% ============================================================================
% Seccion 1: Descripcion del Caso de Uso
% TFP - Rubrica: 1 punto
% ============================================================================
\section{Descripci\'on del caso de uso}

% TODO: Contextualizar el problema de orientacion vocacional en Peru.
% Citar fuentes oficiales (INEI, MINEDU, SUNEDU) sobre oferta educativa,
% tasas de desercion, dificultad de acceso a informacion comparada.

% TODO: Justificar por que IA Generativa + MLOps es adecuado para este caso
% (razonamiento conversacional, recomendaciones personalizadas, reproducibilidad).

% TODO: Definir objetivo de negocio (orientar mejor al estudiante) y objetivo
% tecnico (asistente conversacional hibrido LLM + scoring determinista).

Los estudiantes preuniversitarios en el Per\'u enfrentan una dificultad significativa al momento de elegir una carrera profesional y la instituci\'on de educaci\'on superior donde cursarla. La informaci\'on oficial --- salarios de los egresados, costos de matr\'icula, tasas de admisi\'on, ubicaci\'on de sedes, modalidades de ingreso --- se encuentra dispersa en m\'ultiples portales institucionales y es dif\'icil de comparar de manera sistem\'atica.

[Motivacion del proyecto: por que un copiloto conversacional resuelve este problema mejor que las evaluaciones vocacionales estaticas tradicionales.]

\subsection{Problem\'atica del negocio}

[Describir la problematica concreta con datos. Ejemplos:
- Fragmentacion de la informacion oficial (SUNEDU, portales universitarios).
- Evaluaciones vocacionales estaticas que no consideran restricciones reales (presupuesto, geografia).
- Falta de explicabilidad en recomendaciones automatizadas.]

\subsection{Objetivo de negocio}

[Que decision del usuario queremos mejorar. Ej: reducir el tiempo promedio de decision vocacional, aumentar la satisfaccion con la carrera elegida, etc.]

\subsection{Objetivo t\'ecnico}

[Construir un sistema que:
- Interprete lenguaje natural del estudiante (filtros + chat libre).
- Genere pesos din'amicos para un motor de scoring determinista.
- Devuelva un ranking Top-N de combinaciones Carrera-Universidad con justificacion en lenguaje natural.
- Sea reproducible, versionado y desplegado en AWS.]

\subsection{Fuentes de relevancia}

[Listar fuentes que respalden la relevancia del caso. Ej: estudios de desercion universitaria, reportes de SUNEDU sobre la oferta educativa, papers sobre sistemas de recomendacion en educacion (EdTech).]